\documentclass[unicode,11pt,a4paper,oneside,numbers=endperiod,openany]{scrartcl}

% Required package
\usepackage{amssymb}
\usepackage{graphicx}
\usepackage{amsmath}
\usepackage{matlab-prettifier}
\usepackage{float}
\usepackage[export]{adjustbox}
\usepackage{multirow}
\usepackage{booktabs}
\usepackage{amsthm} % math theorems
\usepackage{ifthen}
\usepackage{physics} % responsive norm, abs, ...
\usepackage{algorithm}
\usepackage{algpseudocode}
\usepackage{array} % for table
\usepackage{mathtools} % for \coloneq
\usepackage{xcolor, soul} % highlight
\usepackage{enumitem} % custom enumeration in lists
\usepackage{dsfont} % 1 for identity matrix
\usepackage{blkarray, bigstrut} % block matrix with labels
\usepackage{tikz} % for drawing

\newcommand\topstrut[1][1.2ex]{\setlength\bigstrutjot{#1}{\bigstrut[t]}}
\newcommand\botstrut[1][0.9ex]{\setlength\bigstrutjot{#1}{\bigstrut[b]}}

\usetikzlibrary{positioning}
\tikzset{
    > = stealth,
    every path/.append style = {
        arrows = ->,
    }
}

\renewcommand{\thesubsection}{\arabic{subsection}}

\newtheorem{theorem}{Theorem}[section]
% 1: title, 2: label, 3: content
\newcommand{\myth}[3]{
    \begin{theorem}[#1] 
        \label{#2} 
        #3 
    \end{theorem}
}

% 1: color, 2: content
\newcommand{\mathcolorbox}[2]{\colorbox{#1}{\(\displaystyle #2\)}}

% 1: if numbered equation, 2: label, 3: content
\newcommand{\myex}[3]{
    \ifthenelse{\equal{#1}{true}}{
        \begin{equation} \label{#2} \begin{aligned} #3 \end{aligned} \end{equation}
    }{
        \begin{equation*} \label{#2} \begin{aligned} #3 \end{aligned} \end{equation*}
    }
}

% vector shortcut
\newcommand{\myvec}[1]{\begin{bmatrix} #1 \end{bmatrix}}
\newcommand{\mymat}[1]{\begin{pmatrix} #1 \end{pmatrix}}

% 1: letter of vector, 2: 
\newcommand\vibar[2]{\mathbf{\overline{#1}}_#2}

% 1: caption, 2: label, 3: trim, 4: figure within figures folder
\newcommand{\myfigure}[4]{
    \begin{figure}[htbp]
    \centering
    \caption{#1}
    \label{#2}
    \includegraphics[width=\paperwidth, trim=#3]{./figures/#4}
    \end{figure}
}

% proof step
\newcommand{\pstep}{\overset{.}{\Longrightarrow}}

\def\ex2f{f(\mathbf{x}^*)}

\usepackage{ifthen}
\usepackage[utf8]{inputenc}
\usepackage{graphics}
\usepackage{graphicx}
\usepackage{hyperref}
\usepackage{amsmath}

\pagestyle{plain}
\voffset -5mm
\oddsidemargin  0mm
\evensidemargin -11mm
\marginparwidth 2cm
\marginparsep 0pt
\topmargin 0mm
\headheight 0pt
\headsep 0pt
\topskip 0pt        
\textheight 255mm
\textwidth 165mm

\newcommand{\duedate} {}
\newcommand{\setduedate}[1]{%
\renewcommand\duedate {\textbf{Due date:}~ #1}}
\newcommand\isassignment {false}
\newcommand{\setassignment}{\renewcommand\isassignment {true}}
\newcommand{\ifassignment}[1]{\ifthenelse{\boolean{\isassignment}}{#1}{}}
\newcommand{\ifnotassignment}[1]{\ifthenelse{\boolean{\isassignment}}{}{#1}}

\newcommand{\assignmentpolicy}{
\begin{table}[h]
\begin{center}
\scalebox{0.8} {%
\begin{tabular}{|p{0.02cm}p{16cm}|}
\hline
&\\
\multicolumn{2}{|c|}{\Large\textbf{Numerical Computing 2023 ---  Submission Instructions}}\\
\multicolumn{2}{|c|}{\large\textbf{(Please, notice that following instructions are mandatory: }}\\
\multicolumn{2}{|c|}{\large\textbf{submissions that don't comply with, won't be considered)}}\\
&\\
\textbullet & Assignments must be submitted to \href{https://www.icorsi.ch/course/view.php?id=14666}{iCorsi} (i.e. in electronic format).\\
\textbullet & Provide both executable package and sources (e.g. C/C++ files, MATLAB). 
If you are using libraries, please add them in the file. Sources must be organized in directories called:\\
\multicolumn{2}{|c|}{\textit{Project\_number\_lastname\_firstname}}\\
& and  the  file must be called:\\
\multicolumn{2}{|c|}{\textit{project\_number\_lastname\_firstname.zip}}\\
\multicolumn{2}{|c|}{\textit{project\_number\_lastname\_firstname.pdf}}\\
\textbullet &  The TAs will grade your project by reviewing your project write-up, and looking at the implementation you attempted, and benchmarking your code's performance.\\

\textbullet & You are allowed to discuss all questions with anyone you like; however: (i) your submission must list anyone you discussed problems with and (ii) you must write up your submission independently.\\
\hline
\end{tabular}
}
\end{center}
\end{table}
}
\newcommand{\punkte}[1]{\hspace{1ex}\emph{\mdseries\hfill(#1~\ifcase#1{Points}\or{Points}\else{Points}\fi)}}


\newcommand\serieheader[6]{
\thispagestyle{empty}%
\begin{flushleft}
\includegraphics[width=0.45\textwidth]{CI_logo}
\end{flushleft}
  \noindent%
  {\large\ignorespaces{\textbf{#1}}\hspace{\fill}\ignorespaces{ \textbf{#2}}}\\ \\%
  {\large\ignorespaces #3 \hspace{\fill}\ignorespaces #4}\\
  \noindent%
  \bigskip
  \hrule\par\bigskip\noindent%
  \bigskip {\ignorespaces {\Large{\textbf{#5}}}
  \hspace{\fill}\ignorespaces \large \ifthenelse{\boolean{\isassignment}}{\duedate}{#6}}
  \hrule\par\bigskip\noindent%  \linebreak
 }

\makeatletter
\def\enumerateMod{\ifnum \@enumdepth >3 \@toodeep\else
      \advance\@enumdepth \@ne
      \edef\@enumctr{enum\romannumeral\the\@enumdepth}\list
      {\csname label\@enumctr\endcsname}{\usecounter
        {\@enumctr}%%%? the following differs from "enumerate"
	\topsep0pt%
	\partopsep0pt%
	\itemsep0pt%
	\def\makelabel##1{\hss\llap{##1}}}\fi}
\let\endenumerateMod =\endlist
\makeatother




\usepackage{textcomp}





\begin{document}

\setassignment
\setduedate{Thursday, 24 October 2024, 23:59}

\serieheader
{Discrete Mathematics}
{2024}
{%
\textbf{Student:} Jeferson Morales Mariciano 
\href{mailto:jmorale@ethz.ch}{\(<\)jmorale@ethz.ch\(>\)} \\\\}
{\vspace{-1cm}}%\textbf{Discussed with:} }
{Assignment 5}{}

%----------------------------------------------------------------------------------
\section*{Exercise 5.5, Properties of Relations (\(\star\)) \hfill (8 Points)}

Prove or disprove the following claims:

\begin{enumerate}[label=\textbf{\alph*)}]
    \item A relation \( \rho \) on a set \( A \) is symmetric on \( A \) 
        if and only if \( \rho^2 \) is symmetric on \( A \).

    \item If \( \rho \) is a relation on a set \( A \) 
        that is symmetric and antisymmetric,
        then it must hold \( \rho = \textsf{id}_A \).

    \item Define the relations \( \rho_1 \) and \( \rho_2 \) on \( \mathbb{Z} \) as
    \[
        a \, \rho_1 \, b \iff b = a + 1, \qquad a \, \rho_2 \, b \iff b \equiv_2 a.
    \]
    Then for \( \rho = \rho_1 \cup \rho_2 \) it holds \( \rho^2 = \mathbb{Z} \times \mathbb{Z} \).
\end{enumerate}


\subsection*{a)}

The claim is false, a counterexample follows:

\myex{false}{ex5-5-a-counterexample}{
    \rho = \{ (a, c), (b, d), (c, b), (d, a) \}
    \\
    \rho^2 = \{ (a, b), (b, a), (c, d), (d, c) \}
}

The relation \( \rho^2 \) is symmetric, but \( \rho \) is not, 
which disprove the claim if and only if (\( \iff \)) from the right to left part 
(\( \Longleftarrow \)). 

The intuition behind the counterexample is given by expanding the right hand side (\( \Longleftarrow \))  
of the claim into a composition of implications as follows. 
Note the universe \( \mathcal{U} \) is an arbitrary set \( A \) of elements.

\myex{false}{ex5-5-a-proof}{
    & \rho^2 \text{ symmetric on } A 
    \\ & \pstep \; 
    \forall a \forall b \left( a \, \rho^2 \, b \Longleftrightarrow b \, \rho^2 \, a \right)
    \hspace{1cm} & \text{(def 3.15 symmetry)}
    \\ & \pstep \; 
    \forall a \forall b \left( (a, b) \in \rho^2 \land (b,a) \in \rho^2 \right)
    \hspace{1cm} & \text{(def 3.9 relation)}
    \\ & \pstep \; 
    \forall a \forall b \left( 
        \exists c \left( (a, c) \in \rho \land (c,b) \in \rho \right)
        \land
        \exists d \left( (b, d) \in \rho \land (d,a) \in \rho \right)
    \right)
    \hspace{1cm} & \text{(def 3.12 composition)}
}

So it is sufficient to find a relation \( \rho \) that is not symmetric, 
i.e. which have \( c \neq d \), satisfying the constraint of the composition of implications above. 
A visualization is provided in the following matrix:

\myex{false}{ex5-5-a-matrix}{
    M^{\rho^2} = 
    \begin{blockarray}{l c c c c}
        & a & b & c & d \\
    \begin{block}{l [c c c c]}
        a & 0 & 1 & 0 & 0 \topstrut \\
        b & 1 & 0 & 0 & 0 \\
        c & 0 & 0 & 0 & 1 \\
        d & 0 & 0 & 1 & 0 \botstrut \\
    \end{block}
    \end{blockarray}
    =
    M^{\rho} \cdot M^{\rho}
    \,,\quad \text{ where } \quad
    M^{\rho} =
    \begin{blockarray}{l c c c c}
        & a & b & c & d \\
    \begin{block}{l [c c c c]}
        a & 0 & 0 & 1 & 0 \topstrut \\
        b & 0 & 0 & 0 & 1 \\
        c & 0 & 1 & 0 & 0 \\
        d & 1 & 0 & 0 & 0 \botstrut \\
    \end{block}
    \end{blockarray}
}

The matrices resemble the counterexample, where \( \rho^2 \) is symmetric,
while \( \rho \) is clearly not, 
i.e. \( M^{\rho} \neq (M^{\rho})^{\intercal} \iff \rho \neq \hat{\rho} \).

A graph representation is provided for the sake of visualization, 
where the composition of the same structure choosen by the matrices is visualize.
The elements of \( \rho^2 \) are given by choosing a path with length 2,
meaning a green edge and the subsequent black one on the right.
Doing so for every vertex gives the \( \rho^2 \) relation. \\

% left graph for rho
\begin{minipage}{0.4\textwidth}
\centering
graph \( \rho \) \\
\vspace{0.2cm}
\tikz{
    \node (a) {$a$};
    \node (b) [below = of a] {$b$};
    \node (c) [below = of b] {$c$};
    \node (d) [below = of c] {$d$};

    \path (a) edge [bend left=10, in=240] (c);
    \path (c) edge (b);
    \path (b) edge [bend left=30] (d);
    \path (d) edge [bend left=20, in=135] (a);
}
\end{minipage}
%
% right graph for rho composition
\begin{minipage}{0.4\textwidth}
\centering
graph of \( \rho^2 \) \\
\vspace{0.2cm}
\tikz{
    \node (al) {$a$};
    \node (bl) [below = of al] {$b$};
    \node (cl) [below = of bl] {$c$};
    \node (dl) [below = of cl] {$d$};
    \node (ar) [right = of al] {$a$};
    \node (br) [below = of ar] {$b$};
    \node (cr) [below = of br] {$c$};
    \node (dr) [below = of cr] {$d$};

    \path (al) edge [green!70!black] (cr);
    \path (bl) edge [green!70!black] (dr);
    \path (cl) edge [green!70!black] (br);
    \path (dl) edge [green!70!black] (ar);

    \path (ar) edge [bend left=10, in=240] (cr);
    \path (cr) edge (br);
    \path (br) edge [bend left=30] (dr);
    \path (dr) edge [bend left=20, in=135] (ar);
}
\end{minipage}


%-----------------------------------------------------------------------------------
\subsection*{b)}

The claim is false, a counterexample follows:

\myex{false}{ex5-5-b-counterexample}{
    \rho = \emptyset \neq \textsf{id}_A
}

The relation \( \rho \) is both symmetric and antisymmetric.
However, \( \rho \) is not the identity relation as the claim states.

The satisfiability conditions for the relation 
being both symmetric and antisymmetric are explained as follows.
The matrix for the counterexample is the all-zeros matrix \( \mathbf{0} \) given by:

\myex{false}{ex5-5-b-matrix}{
    M^{\rho} = 
    \mathbf{0} =
    \myvec{
        0       & \cdots    & 0         \\
        \vdots  & 0         & \vdots    \\
        0       & \cdots    & 0
    }
}

It is clear that the relation is symmetric by Definition 3.15,
\( \forall a \forall b \in A \;\; a \, \rho \, b \iff b \, \rho \, a \),
i.e. the matrix is symmetric 
\( M^{\rho} = (M^{\rho})^{\intercal} \iff \rho = \hat{\rho} \).

Moreover, the relation is antisymmetric by Definition 3.16,
\( \forall a \forall b \in A \;\; a \, \rho \, b \, \land \, b \, \rho \, a \implies a = b \),
i.e. the intersection between the relation and its inverse is a subset of the indentity relation,
\( \rho \cap \hat{\rho} \subseteq \textsf{id}_A \).
Notice that for the counterexample, the intersection is the empty set,
which is a subset of every set by Lemma 3.6, including the identity relation.

\myex{false}{ex5-5-b-antisymmetric}{
    \rho = \emptyset \subseteq \left( A \times A \right),
    \quad 
    \rho \cap \hat{\rho} = \emptyset \subseteq \textsf{id}_{A}
    \quad \because \quad 
    \emptyset \in \mathcal{P(\mathsf{id}_A)}
}

In general symmetry and antisymmetry are independent properties of a relation.

%-----------------------------------------------------------------------------------
\subsection*{c)}

The claim is true, and it can be proven as follows: \\
\noindent recall that the universe \( \mathcal{U} = \mathbb{Z} \), 
so relations will be a subset of the cartesian product \( \mathcal{U} \times \mathcal{U} \),
i.e. \( \rho \subseteq \mathbb{Z} \times \mathbb{Z} \).

\noindent By analysing both relations: 
\begin{itemize}
\item \( \rho_1 = \{ (a, b) \mid b = a + 1 \} \), 
    meaning that the relation denotes the successor \( b \) of \( a \),
    as usually denoted by \( s(\cdot) \). 
    It's matrix representation is scratched hereby:

    \myex{false}{ex5-5-c-rho1-matrix}{
        M^{\rho_1} = 
        \begin{blockarray}{l c c c c c c}
            & \cdots & -1 & 0 & 1 & 2 & \cdots \\
        \begin{block}{l [c c c c c c]}
            \vdots  & \ddots    & \vdots    & \vdots    & \vdots    & \vdots    & \cdots \topstrut \\
            -2      & \cdots    & 1         & 0         & 0         & 0         & \cdots \\
            -1      & \cdots    & 0         & 1         & 0         & 0         & \cdots \\
            0       & \cdots    & 0         & 0         & 1         & 0         & \cdots \\
            1       & \cdots    & 0         & 0         & 0         & 1         & \cdots \\
            \vdots  & \vdots    & \vdots    & \vdots    & \vdots    & \vdots    & \ddots \botstrut \\
        \end{block}
        \end{blockarray}
    }

\item \( \rho_2 = \{ (a, b) \mid b \equiv_2 a \} \), 
    meaning that the relation denotes all elements with same parity, 
    including the same element, thus reflexive since \( \mathsf{id}_{\mathbb{Z}} \subseteq \rho_2 \).
    It is symmetric on \( \mathbb{Z} \) from Example 3.25;
    and it is also transitive from same parity to same parity as stated in Example 3.28.
    The definition of the \( \equiv_m \) relation on \( \mathbb{Z} \) from Example 3.11 
    states \( a \equiv_m b \overset{def}{\iff} \exists k \left( a - b = km \right) \).
    It's matrix representation is scratched hereby:

    \myex{false}{ex5-5-c-rho2-matrix}{
        M^{\rho_2} = 
        \begin{blockarray}{l c c c c c c}
            & \cdots & even & odd & even & odd & \cdots \\
        \begin{block}{l [c c c c c c]}
            \vdots  & \ddots    & \ddots     & \ddots & \ddots & \ddots    & \ddots \topstrut \\
            even    & \ddots         & 1     & 0  & 1  & 0 &  \ddots \\
            odd     & \ddots         & 0     & 1  & 0  & 1 & \ddots \\
            even    & \ddots         & 1     & 0  & 1  & 0 & \ddots \\
            odd     & \ddots         & 0     & 1  & 0  & 1 & \ddots \\
            \vdots  & \ddots         & \ddots     & \ddots & \ddots & \ddots     & \ddots \botstrut \\
        \end{block}
        \end{blockarray}
    }

\end{itemize} 

The informal overview of the proof is: decomposing the claim through composition of implications 
we will try a direct proof
which will lead to a clear satisfiability statement. 
Then, proving by case distinction that the statement is always satisfiable 
w.r.t. fixed interpretation of universe, meaning of predicates and relation symbols.

\subsubsection*{Composition of Implications}

from the statement \( S \) about \textbf{definitions of \( \rho_1, \rho_2 \) with 
\( \rho = \rho_1 \cup \rho_2 \) it holds \( \rho^2 = \mathbb{Z} \times \mathbb{Z} \) }. 
Assuming \( S \) is true, then

\myex{false}{ex5-5-c-implications}{
    \rho^2 
    & \pstep \;
    \forall a \, \forall b \; \left( 
        (a, b) \in \rho^2 
    \right)
    \hspace{1cm} & \text{(def 3.9 relations)}
    \\ & \pstep \; 
    \forall a \, \forall b \; \left( 
        (a, b) \in \left( \rho \circ \rho \right) 
    \right)
    \hspace{1cm} & \text{(def 3.12 composition)}
    \\ & \pstep \; 
    \forall a \, \forall b \, \exists c \; \left( 
        (a, c) \in \rho \land (c, b) \in \rho
    \right)
    \hspace{1cm} & \text{(def 3.12 composition)}
    \\ & \pstep \; 
    \forall a \, \forall b \, \exists c \; \left( 
        (a, c) \in \left( \rho_1 \cup \rho_2 \right) 
        \land 
        (c, b) \in \left( \rho_1 \cup \rho_2 \right) 
    \right)
    \hspace{1cm} & \text{(def of \( \rho \))}
    \\ & \pstep \; 
    \forall a \, \forall b \, \exists c \; \left( 
        \left( 
            (a, c) \in \rho_1       
            \lor
            (a, c) \in \rho_2
        \right) 
        \land 
        \left( 
            (c, b) \in \rho_1
            \lor
            (c, b) \in \rho_2
        \right)
    \right)
    \hspace{1cm} & \text{(def of \( \cup \))}
    \\ & \pstep \; 
    \forall a \, \forall b \, \exists c \; \left( 
        \left( 
            \left( c = a + 1 \right)        
            \lor
            \left( c \equiv_2 a \right)        
        \right) 
        \land 
        \left( 
            \left( b = c + 1 \right)        
            \lor
            \left( b \equiv_2 c \right)        
        \right)
    \right)
    \hspace{1cm} & \text{(def of \( \cup \))}
}

Check that the statement is always satisfiable,
i.e. that there always exists a suitable \( \exists c \in \mathbb{Z} \)
to satisfy the claim.

\subsubsection*{Case Distinction}
For all \( a, b \) in \( \mathbb{Z} \), 
\( a \) can be either even or odd, and same for \( b \).
Recall that the sum of 2 even numbers is even, and also the sum of 2 odd numbers is even in \( \mathbb{Z} \).
The possible cases are then 4, 
and are represented and solved by a suitable \( c \) by the following matrix:

\[
    \begin{array}{|cc|c|}
        \hline
        a & b & c \\
        \hline
        even & even & even    \\
        even & odd & s(a) \\
        odd & even & s(a) \\
        odd & odd & odd    \\
        \hline
    \end{array}
\]

where \( s(x) \) is the successor of \( x \) in the universe.

\noindent Explanation:
\begin{itemize}
\item for \( a \) even and \( b \) even, \( c \) is even: 
    \( (a, c \text{ even}) \circ (c \text{ even}, b) \)   

\item for \( a \) even and \( b \) odd, \( c \) is successor of \( a \): 
    \( (a, a + 1) \circ (a + 1, b) \)   

\item for \( a \) odd and \( b \) even, \( c \) is successor of \( a \):
    \( (a, a + 1) \circ (a + 1, b) \)   

\item for \( a \) odd and \( b \) odd, \( c \) is odd:  
    \( (a, c \text{ odd}) \circ (c \text{ odd}, b) \)   

\end{itemize}

Finally, it follows that the claim is always satisfiable because there always exists a suitable \( c \) 
for the statement, 
meaning is true for all cases of \( a, b \in \mathcal{U} = \mathbb{Z} \times \mathbb{Z} \).
Concluding, \( \rho^2 = \mathds{1} = \mathbb{Z} \times \mathbb{Z} \).


\end{document}
